\section{Capítulo 8 --- Revisão integrada da Unidade I}

Uma cooperativa de entregas decide como reduzir atrasos sem transferir risco injustamente aos entregadores. O caso combina qualidade, estatística, arquitetura, gráficos e comunicação.

\subsection{Questões objetivas}

\ItemHeader{1}{Objetiva}{integração}{Situação profissional e suporte quantitativo}

\TextoBase

O painel declara uma linha por pedido, mas junta avaliações registradas por tentativa. Pedidos com várias tentativas repetem valor e distância, inflando receita e volume. A cooperativa precisa restaurar o grão antes de calcular qualquer indicador de atraso ou produtividade. Uma auditoria encontrou pedidos com até quatro tentativas e confirmou que receita pertence ao cabeçalho, enquanto avaliação pertence à tentativa. A correção precisa preservar as avaliações para diagnóstico sem multiplicar medidas do pedido na camada usada pelo painel. O fechamento financeiro usa o valor do pedido uma única vez, enquanto a operação precisa contar tentativas separadamente. O modelo de consumo deverá oferecer essas duas medidas em grãos explícitos e impedir que um join posterior recombine cabeçalho e detalhe sem agregação prévia ou teste de cardinalidade.

\FonteDidatica

\Comando Selecione a ação que corrige a cardinalidade antes de agregar valores por pedido.

\begin{enumerate}[label=\Alph*)]

\item dividir todos os valores por dois

\item usar DISTINCT no dashboard sem teste

\item corrigir cardinalidade antes de calcular e documentar grão na camada de consumo

\item trocar barras por linhas

\item aceitar porque avaliações são legítimas

\end{enumerate}

\ItemHeader{2}{Objetiva}{centro}{Situação profissional e suporte quantitativo}

\TextoBase

O registro de incidente apresenta no Quadro 1 os sete tempos de entrega do período. O valor 80 foi confirmado após pane e pertence à população observada; a operação quer representar o centro sem ocultar a cauda nem excluir o caso legítimo.
\begin{DocumentBox}[{Quadro 1 --- tempos ordenados, em minutos}]
\centering\begin{tabular}{@{}rrrrrrr@{}}\toprule
18 & 19 & 20 & 21 & 22 & 23 & 80\\\bottomrule
\end{tabular}
\end{DocumentBox}

\FonteDidatica

\Comando Calcule média e mediana e selecione a medida que melhor representa o centro típico, comunicando o extremo.

\begin{enumerate}[label=\Alph*)]

\item mediana 21 representa o centro típico melhor que média 29, pois há cauda extrema

\item média 21

\item moda 20

\item mediana 29

\item amplitude 21

\end{enumerate}

\ItemHeader{3}{Objetiva}{arquitetura}{Situação profissional e suporte quantitativo}

\TextoBase

Eventos brutos precisam permanecer reprocessáveis porque regras de atraso e distância podem ser corrigidas. A cooperativa deseja fontes preservadas, transformação validada e indicadores estáveis por consumidor, com linhagem suficiente para refazer uma publicação contestada.

\FonteDidatica

\Comando Selecione a arquitetura que preserva origem, permite correção e publica indicadores governados com linhagem.

\begin{enumerate}[label=\Alph*)]

\item corrigir o bruto e apagar original

\item calcular tudo no gráfico

\item usar SPEC como nova fonte bruta

\item preservar bruto/bronze ou SOR, transformar em prata/SOT e publicar ouro/SPEC com linhagem

\item duplicar regras em cada painel

\end{enumerate}

\ItemHeader{4}{Objetiva}{probabilidade}{Situação profissional e suporte quantitativo}

\TextoBase

Em 500 entregas comparáveis, 75 atrasaram. Entre 100 realizadas sob chuva, 30 atrasaram. A equipe precisa calcular probabilidade geral e condicional, reconhecendo que chuva pode estar associada a região, horário e distância e que a tabela isolada não demonstra causa.

\FonteDidatica

\Comando Calcule $P(A)$ e $P(A\mid C)$ e selecione a interpretação proporcional à evidência observacional.

\begin{enumerate}[label=\Alph*)]

\item 15\% e 6\%

\item P(atraso)=15\% e P(atraso|chuva)=30\%; há associação, não causalidade isolada

\item 30\% e 15\%

\item chuva causa exatamente metade dos atrasos

\item P(chuva|atraso)=30\%

\end{enumerate}

\ItemHeader{5}{Objetiva}{visual}{Situação profissional e suporte quantitativo}

\TextoBase

Um mapa de calor compara bairros, mas não apresenta escala, usa cores sem ordenação e mistura contagem de atrasos com percentual de entregas. Bairros possuem volumes muito diferentes. A revisão deve preservar unidade, denominador e acessibilidade antes de destacar uma região.

\FonteDidatica

\Comando Selecione o redesenho que separa medidas, explicita escala e permanece legível sem depender apenas de cor.

\begin{enumerate}[label=\Alph*)]

\item aumentar o número de cores

\item retirar legenda

\item somar contagem e percentual

\item ordenar bairros alfabeticamente apenas

\item separar medidas, explicitar escala/unidade e oferecer rótulos acessíveis

\end{enumerate}

\ItemHeader{6}{Objetiva}{IQR}{Situação profissional e suporte quantitativo}

\TextoBase

Para atrasos, $Q_1=8$ e $Q_3=20$ minutos; uma entrega levou 42. A equipe adotou o critério $1{,}5\,IQR$ apenas como sinalizador e confirmou a observação na fonte. Precisa calcular a cerca superior e decidir se o caso merece investigação.

\FonteDidatica

\Comando Calcule o IQR e a cerca superior e selecione o tratamento exploratório apropriado para 42.

\begin{enumerate}[label=\Alph*)]

\item IQR=28 e 42 é comum

\item cerca=32

\item IQR=12, cerca superior=38; 42 é potencial outlier a investigar

\item qualquer valor acima de Q3 deve ser removido

\item outlier prova erro de coleta

\end{enumerate}

\ItemHeader{7}{Objetiva}{semântica}{Situação profissional e suporte quantitativo}

\TextoBase

A expressão ``entrega concluída'' inclui devolução em uma fonte e a exclui em outra. Ambas passam em testes de tipo e completude, mas alimentam o mesmo indicador executivo. Antes de combinar dados, a cooperativa precisa governar definição, finalidade e vigência.

\FonteDidatica

\Comando Selecione a ação de governança necessária antes de combinar ou comparar o indicador entre fontes.

\begin{enumerate}[label=\Alph*)]

\item calcular média das fontes

\item governar definição antes de combinar e nomear métricas distintas se usos diferirem

\item deixar o gráfico resolver

\item usar definição com maior valor

\item ocultar a divergência

\end{enumerate}

\ItemHeader{8}{Objetiva}{correlação}{Situação profissional e suporte quantitativo}

\TextoBase

Distância e atraso apresentam correlação agregada $r=0,70$. Ao segmentar por região, as correlações variam de 0,10 a 0,25, e as regiões possuem perfis de tráfego diferentes. A equipe precisa avaliar efeito de composição antes de afirmar relação operacional forte.

\FonteDidatica

\Comando Compare a correlação agregada às segmentadas e selecione a explicação que considera composição e confundidores.

\begin{enumerate}[label=\Alph*)]

\item 70\% do atraso é causado pela distância

\item segmentação invalida toda correlação

\item somar os coeficientes

\item associação agregada pode refletir composição; analisar segmentos e confundidores

\item usar apenas o maior valor

\end{enumerate}

\ItemHeader{9}{Objetiva}{decisão}{Situação profissional e suporte quantitativo}

\TextoBase

Reduzir a janela prometida elevou o indicador de pontualidade, mas também aumentou acidentes relatados e jornadas prolongadas. A direção não autorizou compensar segurança por desempenho. A decisão precisa tratar qualidade e risco como restrições conjuntas, não como média.

\FonteDidatica

\Comando Selecione a decisão que usa pontualidade e segurança como critérios não compensatórios e verificáveis.

\begin{enumerate}[label=\Alph*)]

\item formular decisão multicritério com limite de segurança não compensável por média

\item otimizar somente pontualidade

\item tirar média de acidentes e tempo

\item ocultar efeito adverso

\item escolher a métrica mais fácil

\end{enumerate}

\ItemHeader{10}{Objetiva}{auditoria}{Situação profissional e suporte quantitativo}

\TextoBase

Um insight gerado por IA afirma que atrasos cresceram, mas não informa consulta, snapshot, período, filtros, regra de atraso nem versão do modelo. A diretoria deseja agir sobre o resultado. Antes, precisa existir evidência suficiente para reproduzir e contestar a resposta.

\FonteDidatica

\Comando Selecione o pacote mínimo de proveniência necessário antes de aprovar o insight para decisão.

\begin{enumerate}[label=\Alph*)]

\item aceitar se a frase for plausível

\item pedir resposta mais confiante

\item validar só a gramática

\item publicar com aviso genérico

\item tratar como hipótese e exigir consulta, fonte, recorte, cálculo e reprodução

\end{enumerate}

\subsection{Questões discursivas}

\ItemHeader{11}{Discursiva}{Integrar evidências e justificar decisão}{Caso profissional do capítulo}

\TextoBase

Uma cooperativa de entregas precisa explicar atrasos a municípios associados. A base mistura pedidos e ocorrências, duas áreas usam definições distintas de ``entrega no prazo'' e o painel omite bairros com baixa cobertura. A diretoria quer uma recomendação que possa ser reproduzida e contestada.

\FonteDidatica

\Comando Estruture uma análise que conecte pergunta decisória, grão, qualidade, medida estatística, arquitetura, visual, narrativa e proteção aos grupos pouco observados. Indique evidências necessárias antes de recomendar uma ação.

\ItemHeader{12}{Discursiva}{Integrar evidências e justificar decisão}{Caso profissional do capítulo}

\TextoBase

Uma amostra de tempos de entrega contém $18,19,20,21,22,23,80$ minutos. No mesmo período, ocorreram 75 atrasos em 500 entregas; durante chuva, foram 30 atrasos em 100 entregas. A operação quer resumir centro, dispersão e associação com chuva.

\FonteDidatica

\Comando Calcule média, mediana e IQR dos tempos; calcule $P(\text{atraso})$, $P(\text{atraso}\mid\text{chuva})$ e a razão entre essas taxas. Interprete o extremo e explique por que a razão não demonstra causalidade.

\ItemHeader{13}{Discursiva}{Integrar evidências e justificar decisão}{Caso profissional do capítulo}

\TextoBase

Uma junção entre pedidos e ocorrências multiplicou linhas; além disso, comercial define prazo pela promessa ao cliente e logística pelo horário planejado. Registros sem município foram descartados silenciosamente. O total publicado diverge do sistema operacional e não pode ser reconstituído.

\FonteDidatica

\Comando Diagnostique fanout, divergência semântica e ausência. Distribua responsabilidades entre SOR, SOT e SPEC e proponha testes de unicidade, reconciliação, completude e contrato antes da publicação.

\ItemHeader{14}{Discursiva}{Integrar evidências e justificar decisão}{Caso profissional do capítulo}

\TextoBase

Três recomendações disputam orçamento: A reduz dois minutos no tempo médio, mas aumenta acidentes; B reduz atrasos em bairros periféricos com efeito médio menor; C maximiza produtividade, porém depende de dados pessoais sem base autorizada. A direção pede um ranking único.

\FonteDidatica

\Comando Compare A, B e C por segurança, efeito, equidade e auditabilidade. Declare um critério não compensável, justifique eliminações e explique quais evidências faltam para escolher entre alternativas admissíveis.

\ItemHeader{15}{Discursiva}{Integrar evidências e justificar decisão}{Caso profissional do capítulo}

\TextoBase

Um assistente de IA afirmou que atrasos cresceram 18\% e recomendou trocar a transportadora. A resposta não informa período, denominador, consulta, versão da métrica nem tratamento de pedidos cancelados. O gestor deseja aprovar a recomendação ainda hoje.

\FonteDidatica

\Comando Audite o insight e liste as evidências mínimas para reproduzir o percentual, contestar sua interpretação e avaliar a recomendação. Defina uma regra objetiva de aprovação, revisão ou recusa.

\newpage

\subsection{Respostas explicadas das questões pares}

\paragraph{Questão 2 --- resposta A.} A soma é $18+19+20+21+22+23+80=203$, logo a média é $203/7=29$. Ordenados os sete valores, a mediana é o quarto: 21. O caso de 80 puxa a média oito minutos acima da mediana; por isso, 21 representa melhor o centro típico, sem autorizar a exclusão do extremo confirmado. Logo, A.
\[\bar x=29\text{ min},\qquad \widetilde{x}=21\text{ min}\]

\paragraph{Questão 4 --- resposta B.} No período, $75/500=0{,}15=15\%$. Sob chuva, $30/100=0{,}30=30\%$. A taxa condicional é o dobro da geral, evidenciando associação no recorte observado. Isso não isola chuva de rota, turno, distância ou outros fatores; portanto não basta para afirmar causalidade. A alternativa B preserva cálculo e interpretação.

\paragraph{Questão 6 --- resposta C.} Com $Q_1=8$ e $Q_3=20$, $IQR=20-8=12$. A cerca superior de Tukey é $Q_3+1{,}5IQR=20+18=38$. Como $42>38$, o valor é sinalizado como potencial outlier. O boxplot pede investigação de origem e contexto, não exclusão automática; por isso, C.

\paragraph{Questão 8 --- resposta D.} Uma correlação agregada pode mudar quando a composição dos grupos muda, mesmo sem alteração da relação dentro de cada segmento. Antes de agir, a equipe deve estratificar por variáveis relevantes, comparar denominadores e procurar confundidores. A alternativa D é a única que trata o coeficiente como evidência a investigar, não como prova causal.

\paragraph{Questão 10 --- resposta E.} Sem período, denominador, fonte e transformação, os 18\% não são reproduzíveis. O texto da IA deve ser tratado como hipótese até que consulta, snapshot, versão da métrica, filtros e cálculo sejam registrados e reexecutados. Só depois se avalia a recomendação; esse encadeamento corresponde a E.

\paragraph{Questão 12 --- critérios.} Esperam-se média $203/7=29$, mediana 21 e, pela convenção de metades sem a mediana, $Q_1=19$, $Q_3=23$ e $IQR=4$. Também: $75/500=15\%$, $30/100=30\%$ e razão $30/15=2$. A interpretação deve reconhecer o extremo e separar associação de causalidade.

\paragraph{Questão 14 --- critérios.} A resposta deve tornar segurança um requisito não compensável, eliminando A; deve recusar C enquanto não houver base autorizada e governança; e deve avaliar B com evidência estratificada, incerteza e monitoramento. Um ranking puramente aditivo que permita compensar dano grave não atende ao comando.
